\documentclass[11pt]{amsart}

\usepackage[T1]{fontenc}
\usepackage{lmodern}
\usepackage{microtype}
\usepackage{mathtools,amssymb,amsthm}
\usepackage{array}
\usepackage[hidelinks]{hyperref}

\hypersetup{
  pdftitle={An Explicit Negatively Associated Determinantal Point Process with No Symmetric Kernel, Minimal at n=3}
}

\newtheorem{theorem}{Theorem}
\newtheorem{proposition}[theorem]{Proposition}
\newtheorem{lemma}[theorem]{Lemma}
\newtheorem{corollary}[theorem]{Corollary}
\theoremstyle{remark}
\newtheorem{remark}[theorem]{Remark}

\DeclareMathOperator{\DPP}{DPP}
\newcommand{\bbR}{\mathbb R}
\newcommand{\bbC}{\mathbb C}

\title[An explicit negatively associated DPP with no symmetric kernel]
{An Explicit Negatively Associated Determinantal Point Process\\ with No Symmetric Kernel,
Minimal at $n=3$}

\date{}

\subjclass[2020]{60G55, 15A15, 60E15, 15B48}
\keywords{determinantal point process, nonsymmetric kernel, negative association,
principal minors, strongly Rayleigh}

\begin{document}

\begin{abstract}
In the discussion section of \cite{Poinas} Poinas asks: \emph{does a negatively associated
determinantal point process that does not have the same distribution as a $\DPP$ with a
symmetric kernel even exist?} Read as a question about negative association in the standard
sense recalled in Section~\ref{sec:question}, it does, already on three points. The
rational kernel
\[
 W=\begin{pmatrix}
  1/2 & 1/10 & 1/25\\ 1/10 & 1/2 & 1/10\\ 1/4 & 1/10 & 1/2
 \end{pmatrix}
\]
is a $\DPP$ kernel, $\DPP(W)$ is negatively associated -- strictly so for nonconstant
increasing functions on nonempty disjoint supports -- and no real symmetric, indeed no
Hermitian, kernel has its law, because the law invariant $\sigma^{2}-4uvw$ equals
$441/10^{8}>0$ while every real symmetric kernel forces $\sigma^{2}-4uvw=0$ and every
Hermitian kernel forces $\sigma^{2}-4uvw\le0$. No smaller ground set carries such a process:
at $n\le2$ a short case check shows every negatively associated $\DPP$ to be symmetrisable.
Two qualifications belong here and not further down. First, $\DPP(W)$ is \emph{not} strongly
Rayleigh, while the sentence motivating the question in \cite{Poinas} attributes negative
association to symmetric-kernel $\DPP$s by citing \cite{BBL,Lyons}, and the theorem of
\cite{BBL} gives those $\DPP$s the stronger strongly Rayleigh property; nothing here settles
the question with
negative association replaced by that property. Second, non-symmetrisable determinantal laws
are not new: Section~\ref{sec:attribution} records two published constructions, one of them in
\cite{Poinas} itself, whose three-point laws already have positive invariant, and for neither
is negative association -- the other half of the question -- established.
\end{abstract}

\maketitle

\section{The question}\label{sec:question}

Fix $n\ge1$ and $K\in M_{n}(\bbR)$. Following \cite[\S2]{Poinas}, the \emph{determinantal
measure} with kernel $K$ is the set function
\begin{equation}\label{eq:defdpp}
 \mu\bigl(\{X:\ A\subseteq X\}\bigr)=\det K_{A}\qquad(A\subseteq[n]),
\end{equation}
where $K_{A}$ is the principal submatrix on $A$ and $\det K_{\emptyset}=1$. Its atoms are
recovered by M\"obius inversion,
\begin{equation}\label{eq:mobius}
 \mathbb P(X=A)=\sum_{A\subseteq T\subseteq[n]}(-1)^{|T|-|A|}\det K_{T},
\end{equation}
and $K$ is called a \emph{$\DPP$ kernel} exactly when all $2^{n}$ of these numbers are
nonnegative, i.e.\ when $\mu$ is a probability measure; $\DPP(K)$ then denotes the resulting
random subset of $[n]$. For \emph{symmetric} $K$, \cite[\S2]{Poinas} records that this
happens iff every eigenvalue of $K$ lies in $[0,1]$.

A random subset $X\subseteq[n]$ is \emph{negatively associated} (NA) if for all
\emph{disjoint} $A,B\subseteq[n]$ and all increasing $f\colon2^{A}\to\bbR$,
$g\colon2^{B}\to\bbR$,
\begin{equation}\label{eq:na}
 \mathbb E\bigl[f(X\cap A)\,g(X\cap B)\bigr]\le
 \mathbb E\bigl[f(X\cap A)\bigr]\,\mathbb E\bigl[g(X\cap B)\bigr].
\end{equation}
Negative association is not defined in \cite{Poinas}; it is imported there by the citations
\cite{BBL,Lyons}, and \eqref{eq:na} is the standard finite-ground-set form. Every $\DPP$
with a symmetric kernel is strongly Rayleigh \cite{BBL} and hence NA. Two laws
$\DPP(K)$, $\DPP(S)$ coincide iff $\det K_{A}=\det S_{A}$ for \emph{every} $A\subseteq[n]$,
because by \eqref{eq:defdpp} the principal minors are the inclusion probabilities.

\medskip
\noindent\textbf{The question, and where it is.}
The paragraph below is quoted from the subsubsection \emph{Negative association} of the
section \emph{Discussion and open problems} of \cite{Poinas}. The question is unnumbered and
in prose.

\begin{quote}
``All DPPs with symmetric kernels satisfy a strong dependency property called negative
association [Borcea, Lyons] that extends to continuous DPPs with numerous applications
[Ghosh, Poinas17]. Since negative association implies pairwise repulsion then a necessary
condition for a DPP with kernel $K$ to be negatively associated is that
$K_{i,j}K_{j,i}\ge 0$ for any $i,j\in[n]$ which is obviously not satisfied by all generic
DPP kernels. The question naturally arising from this conclusion is what condition on $K$ is
needed for the associated DPP to be negatively associated? In fact, does a negatively
associated DPP that does not have the same distribution as a DPP with symmetric kernel even
exist?''
\end{quote}

\noindent
In the source the four bracketed items are \verb|\cite| commands; apart from those the
paragraph is quoted as it stands. The clause settled below is its final sentence.

It is an existence question about \emph{negative association} as defined in \eqref{eq:na}, and
that is the reading taken here. A caution is owed: the paragraph opens by attributing negative
association to symmetric-kernel $\DPP$s, and those satisfy the stronger strongly Rayleigh
property as recorded above, so a reader may take the question to be about strongly Rayleigh
measures instead. The witness below is \emph{not} strongly Rayleigh
(Section~\ref{sec:open}), so under that stronger reading the question is not settled here. The
\emph{first} question of the quoted paragraph -- ``what condition on $K$ is needed?'' -- is a
different target and is not addressed below.

\section{The witness}\label{sec:witness}

\begin{theorem}\label{thm:main}
Let
\[
 W=\begin{pmatrix}
   1/2 & 1/10 & 1/25\\ 1/10 & 1/2 & 1/10\\ 1/4 & 1/10 & 1/2
  \end{pmatrix}.
\]
Then \emph{(i)} $W$ is a $\DPP$ kernel; \emph{(ii)} $\DPP(W)$ is negatively associated, and
\eqref{eq:na} is \emph{strict} whenever $A,B$ are nonempty and $f,g$ are both nonconstant; and
\emph{(iii)} no real symmetric matrix, and no
Hermitian matrix, has the same principal minors as $W$, so $\DPP(W)$ is the law of no
symmetric-kernel and of no Hermitian-kernel $\DPP$. In particular the final question of
\cite[\S{}Discussion and open problems, ``Negative association'']{Poinas}, read as a question
about \eqref{eq:na}, has the answer \textbf{yes}.
\end{theorem}

\begin{proof}[Proof of (i)]
Every $1\times1$ minor is $1/2$. Every $2\times2$ principal minor is
$1/4-1/100=6/25$, because
$W_{12}W_{21}=W_{13}W_{31}=W_{23}W_{32}=1/100$ (namely
$\tfrac1{10}\tfrac1{10}$, $\tfrac1{25}\tfrac14$, $\tfrac1{10}\tfrac1{10}$). By the six-term
Laplace expansion,
\[
 \det W=\tfrac18-3\cdot\tfrac12\cdot\tfrac1{100}
 +W_{12}W_{23}W_{31}+W_{13}W_{32}W_{21}
 =\tfrac18-\tfrac3{200}+\tfrac1{400}+\tfrac1{2500}=\tfrac{1129}{10000}.
\]
The minors depend only on $|A|$, so the law is exchangeable. By \eqref{eq:mobius} the eight
atoms are
\[
 \begin{array}{r@{\quad}l}
  \mathbb P(X=\emptyset) &=1-3\cdot\tfrac12+3\cdot\tfrac6{25}-\tfrac{1129}{10000}
   =\tfrac{1071}{10000},\\[2pt]
  \mathbb P(X=\{i\}) &=\tfrac12-2\cdot\tfrac6{25}+\tfrac{1129}{10000}
   =\tfrac{1329}{10000},\\[2pt]
  \mathbb P(X=\{i,j\}) &=\tfrac6{25}-\tfrac{1129}{10000}=\tfrac{1271}{10000},\\[2pt]
  \mathbb P(X=[3]) &=\tfrac{1129}{10000}.
 \end{array}
\]
Their sum is $(1071+3\cdot1329+3\cdot1271+1129)/10000=1$ exactly, and all eight are
\emph{strictly} positive, the smallest being $1071/10000$. Hence $\mu$ is a probability
measure and $W$ is a $\DPP$ kernel. (Also $\mathbb P(X=\emptyset)=\det(I-W)\ne0$, so
$L=W(I-W)^{-1}$ exists and $\DPP(W)$ is an $L$-ensemble.)
\end{proof}

\begin{proof}[Proof of (ii)]
Write $D(f,g)=\mathbb E[fg]-\mathbb E[f]\mathbb E[g]$, so that \eqref{eq:na} reads
$D(f,g)\le0$. The defect $D$ is bilinear and vanishes as soon as $f$ or $g$ is constant; in
that case \eqref{eq:na} holds with equality, which is why strictness is asserted only for
nonconstant $f,g$. Every increasing $f$ on a finite Boolean lattice is a \emph{nonnegative}
combination of up-set indicators plus a constant, namely
$f=v_{1}+\sum_{i\ge2}(v_{i}-v_{i-1})\mathbf 1_{\{f\ge v_{i}\}}$ over its distinct values
$v_{1}<\dots<v_{m}$; if $f$ is nonconstant then $m\ge2$, each coefficient
$v_{i}-v_{i-1}$ is strictly positive, and each $\{f\ge v_{i}\}$, $i\ge2$, is a proper nonempty
up-closed family. So \eqref{eq:na} for all increasing $f,g$ is \emph{equivalent} to the finite
list of inequalities $D(\mathbf 1_{U},\mathbf 1_{V})\le0$ over proper nonempty up-closed
$U\subseteq2^{A}$, $V\subseteq2^{B}$, not merely implied by it, and if every one of those is
strict then $D(f,g)<0$ for every pair of nonconstant increasing $f,g$.

At $n=3$ the disjoint pairs are $(\{i\},\{j\})$ and $(\{i\},\{j,k\})$, three of each; the
proper nonempty up-sets number $1$ on a one-element lattice and $4$ on a two-element
lattice, so there are $3\cdot1\cdot1+3\cdot1\cdot4=15$ inequivalent inequalities, $30$ when
both orders of each pair are counted. Exchangeability collapses them to three classes, each
\emph{strict}:
\[
 \begin{array}{lccc}
  \text{class} & \text{LHS} & \text{RHS} & \text{slack}\\ \hline
  \{i\in X\}\ \text{vs}\ \{j\in X\} & 6/25 & 1/4 & 1/100\\
  \{i\in X\}\ \text{vs}\ \{j,k\in X\} & 1129/10000 & 3/25 & 71/10000\\
  \{i\in X\}\ \text{vs}\ \{X\cap\{j,k\}\ne\emptyset\} & 3671/10000 & 19/50 & 129/10000
 \end{array}
\]
All $30$ ordered inequalities hold, with minimum slack $71/10000>0$. (The first class is
exactly the necessary condition $K_{ij}K_{ji}\ge0$ of \cite{Poinas} quoted in
Section~\ref{sec:question}; here all three products equal $1/100$.)
\end{proof}

For (iii) we isolate the invariant.

\section{The obstruction is a single law invariant}\label{sec:invariant}

For a triple $\{i,j,k\}\subseteq[n]$ and $K\in M_{n}(\bbC)$ put
\begin{equation}\label{eq:inv}
\begin{gathered}
 u=K_{ij}K_{ji},\qquad v=K_{ik}K_{ki},\qquad w=K_{jk}K_{kj},\\
 p=K_{ij}K_{jk}K_{ki},\qquad q=K_{ik}K_{kj}K_{ji},\qquad \sigma=p+q.
\end{gathered}
\end{equation}

\begin{lemma}\label{lem:pq}
$pq=uvw$ identically, with no hypothesis on $K$; consequently
$\sigma^{2}-4uvw=(p-q)^{2}$. Moreover $u,v,w$ and $\sigma$ are functions of the principal
minors of $K$ alone:
\[
 u=K_{ii}K_{jj}-\det K_{\{i,j\}},\qquad
 \sigma=\det K_{\{i,j,k\}}-K_{ii}K_{jj}K_{kk}+K_{ii}w+K_{jj}v+K_{kk}u .
\]
\end{lemma}

\begin{proof}
$pq=(K_{ij}K_{ji})(K_{jk}K_{kj})(K_{ki}K_{ik})=uvw$ after regrouping the six factors, and
$(p+q)^{2}-4pq=(p-q)^{2}$. The first displayed identity is the $2\times2$ expansion; the
second is the six-term expansion of $\det K_{\{i,j,k\}}$ with $p+q$ collected.
\end{proof}

\begin{theorem}[representability]\label{thm:A}
Fix $u,v,w>0$ and a value $\sigma$. Among matrices with the invariants \eqref{eq:inv} on a
given triple:
a real matrix exists iff $\sigma^{2}\ge4uvw$; a real \textbf{symmetric} one exists iff
$\sigma^{2}=4uvw$; a \textbf{Hermitian} one exists iff $\sigma^{2}\le4uvw$.
\end{theorem}

\begin{proof}
By Lemma~\ref{lem:pq}, $p$ and $q$ are the two roots of $z^{2}-\sigma z+uvw$. For real
matrices they must be real, i.e.\ $\sigma^{2}\ge4uvw$. For real \emph{symmetric} $S$ the two
$3$-cycle products coincide, $p=q=S_{ij}S_{jk}S_{ki}$, so $\sigma=2p$ and $uvw=p^{2}$, i.e.\
$\sigma^{2}=4uvw$; conversely $\sigma^{2}=4uvw$ with $u,v,w>0$ is realised by
$S_{ij}=\varepsilon_{ij}\sqrt u$ etc.\ with signs chosen so that the product is
$\sigma/2$. For Hermitian $S$ one has $q=\bar p$, so $\sigma=2\operatorname{Re}p$ and
$uvw=|p|^{2}$, whence
\begin{equation}\label{eq:herm}
 \sigma^{2}-4uvw=4(\operatorname{Re}p)^{2}-4|p|^{2}=-4(\operatorname{Im}p)^{2}\le0;
\end{equation}
conversely any $\sigma$ with $\sigma^{2}\le4uvw$ is realised by a $p$ on the circle
$|p|^{2}=uvw$ with $2\operatorname{Re}p=\sigma$.
\end{proof}

\begin{proof}[Proof of Theorem~\ref{thm:main}(iii)]
For $W$ on $\{1,2,3\}$: $u=v=w=1/100$, $p=W_{12}W_{23}W_{31}=1/400$,
$q=W_{13}W_{32}W_{21}=1/2500$, so $\sigma=29/10000$ and, by Lemma~\ref{lem:pq},
\[
 \sigma^{2}-4uvw=(p-q)^{2}=\Bigl(\tfrac1{400}-\tfrac1{2500}\Bigr)^{2}
 =\Bigl(\tfrac{21}{10000}\Bigr)^{2}=\frac{441}{10^{8}}>0 .
\]
Since $u,v,w,\sigma$ are functions of the principal minors alone, any matrix with the law of
$\DPP(W)$ has the same value $441/10^{8}$; by Theorem~\ref{thm:A} that excludes every real
symmetric matrix (which needs $0$) and, by \eqref{eq:herm}, every Hermitian matrix (which
needs $\le0$).
\end{proof}

\begin{remark}[an algebra-free version of the same exclusion]\label{rem:signs}
A real symmetric $S$ with the principal minors of $W$ must have $S_{ii}=1/2$ and
$S_{ij}^{2}=W_{ii}W_{jj}-\det W_{\{i,j\}}=1/100$, hence $S_{ij}=\pm1/10$: the modulus is
\emph{forced}, and $1/100$ is a rational square, so no irrationality escapes. Enumerating the
eight sign patterns gives only two values of $\det S$, namely $14/125$ and $27/250$, and
neither is $\det W=1129/10000$. This settles (iii) for the real symmetric case with no
algebra at all.
\end{remark}

\section{Minimality in $n$}\label{sec:sharp}

\begin{theorem}\label{thm:D}
A negatively associated $\DPP$ whose law is the law of no symmetric-kernel $\DPP$ exists for
\textbf{exactly} the $n\ge3$. In particular $n=3$ is the smallest possible ground set.
\end{theorem}

\begin{proof}
\emph{$n\ge3$.} Take $K_{n}=W\oplus\operatorname{diag}(c_{4},\dots,c_{n})$ with rationals
$c_{i}\in(0,1)$. For block-diagonal $K$ the minors factorise over blocks -- this is
\cite[\S2]{Poinas}: ``if $K$ is a block diagonal matrix with diagonal blocks
$K_{1},\dots,K_{l}$ then $K$ being a $\DPP$ kernel is equivalent to each $K_{i}$ being a
$\DPP$ kernel'' -- so the atoms factorise, all stay strictly positive, and $\DPP(K_{n})$ is
the independent superposition of $\DPP(W)$ with independent Bernoulli$(c_{i})$'s. Negative
association is closed under independent superposition \cite[property P7]{JDP}, so
$\DPP(K_{n})$ is NA. The obstruction is inherited by the principal submatrix on
$\{1,2,3\}$: any symmetric or Hermitian $L$ with the same law has the same principal minors,
hence the same $u,v,w,\sigma$ on $\{1,2,3\}$, hence $\sigma^{2}-4uvw\le0$ there, against
$441/10^{8}>0$. No information about $L$ outside $\{1,2,3\}$ is used, so one bad triple
obstructs at every larger $n$.

\emph{$n\le2$ is impossible.} At $n=1$, \eqref{eq:na} is vacuous (there is no pair of
disjoint nonempty subsets) and $K=(K_{11})$ is symmetric anyway. At $n=2$ the law is
determined by $(K_{11},K_{22},u)$ with $u=K_{12}K_{21}$, and the single NA inequality is
$\mathbb P(1,2\in X)\le\mathbb P(1\in X)\mathbb P(2\in X)$, i.e.\
$K_{11}K_{22}-u\le K_{11}K_{22}$, i.e.\ $u\ge0$. Then the symmetric
$S=\begin{psmallmatrix}K_{11}&\sqrt u\\ \sqrt u&K_{22}\end{psmallmatrix}$ reproduces both
$1\times1$ minors and the $2\times2$ minor $K_{11}K_{22}-u$, hence every atom; and $S$ has
spectrum in $[0,1]$ because its characteristic polynomial satisfies
$\chi(0)=\mathbb P(X=\{1,2\})\ge0$, $\chi(1)=\mathbb P(X=\emptyset)\ge0$ and
$0\le\tfrac12\operatorname{tr}S\le1$. So every NA $\DPP$ on at most two points is
symmetrisable.
\end{proof}

\section{Real-symmetric-kernel laws lie on a hypersurface}\label{sec:generic}

Coordinatise the $3$-point $\DPP$ laws by the $7$ invariants
$(a_{1},a_{2},a_{3};u,v,w;\sigma)$ with $a_{i}=K_{ii}$; by Lemma~\ref{lem:pq} these are exactly
the principal minors, so they are exactly the law.

\begin{theorem}\label{thm:generic}
The negatively associated $3$-point $\DPP$ laws have nonempty interior in this
$7$-dimensional space, whereas the laws of real symmetric kernels are confined to the
hypersurface $\sigma^{2}=4uvw$, a closed Lebesgue-null set of codimension one. Consequently
\textbf{Lebesgue-almost every negatively associated $3$-point $\DPP$ law is the law of no
$\DPP$ with a real symmetric kernel.}
\end{theorem}

The restriction to real symmetric kernels is essential: Hermitian kernels are confined only to
the region $\sigma^{2}\le4uvw$, which is full-dimensional, so no analogous genericity statement
against Hermitian kernels is made here; see Section~\ref{sec:open}.

\begin{proof}
At the law of $W$ all eight atom inequalities, all $30$ NA inequalities, $u,v,w>0$ and
$\sigma^{2}-4uvw>0$ hold \emph{strictly} (Section~\ref{sec:witness}), and each is a
polynomial in the seven free coordinates; so an open neighbourhood of that law consists of NA
$\DPP$ laws with $\sigma^{2}>4uvw$. The symmetric locus is $\{\sigma^{2}=4uvw\}$ by
Theorem~\ref{thm:A}, the zero set of a nonzero polynomial.
\end{proof}

\section{What was already in print}\label{sec:attribution}

\emph{Determinantal laws that are the law of no symmetric kernel are already in print, in two
published constructions; whether either is negatively associated we do not know.} We record
this because presenting non-symmetrisability as a discovery would be an overclaim.

\medskip
\noindent\textbf{(a) The asking paper's own coupling.} Poinas' coupling construction
\cite[Prop.~``Construction coupling'']{Poinas} builds
from a symmetric $\DPP$ kernel $K$ and a parameter vector $p$ the $2n$-site kernel
$\mathcal K=\begin{psmallmatrix}K&K-D(p)\\K&K\end{psmallmatrix}$. On the triple
$\{1,2,1'\}$ (two sites of the first copy, one of the second) one gets $u=w=b^{2}$ with
$b=K_{12}$, $v=K_{11}(K_{11}-p_{1})$ and $\sigma^{2}-4uvw=(b^{2}p_{1})^{2}>0$ whenever
$p_{1}>0$ and $b\ne0$. So the paper's own coupling law is already non-symmetrisable; should it
also be negatively associated -- which we have not checked -- bare existence would follow from
that construction alone. Poinas does not test it for negative association.

\medskip
\noindent\textbf{(b) The one-dependent class the paper itself cites.} Borodin, Diaconis and
Fulman \cite{BDF} -- reference \texttt{[Borodin09]} of \cite{Poinas} -- construct
determinantal two-block-factor processes with nonsymmetric kernels; their
$X_{i}=\mathbf 1\{Y_{i}\le p,\,Y_{i+1}>p\}$ gives, at $n=3$,
$K=d\begin{psmallmatrix}1&1&0\\1&1&1\\1&1&1\end{psmallmatrix}$ with $d=p(1-p)$, whose law
forces $\det K=0$ while a symmetric competitor would need $-d^{3}$. Their Theorem~4.1 and
Example~6 give $\mathbb P(0,0,0)=1-3d+d^{2}$ at $a_{2}=d$, $a_{3}=a_{4}=0$ -- the same atom we
obtain by M\"obius inversion from that kernel, an external published checksum. Not every
member works: their descent process of a uniform permutation on three sites has $u=w=1/12$,
$v=0$, $\sigma=0$, hence $\sigma^{2}=4uvw$, and \emph{is} symmetrisable.

\medskip
\noindent\textbf{What this leaves.} In the sources we read we did not find negative
association stated for any nonsymmetric $\DPP$ kernel, nor the question whether an NA $\DPP$
law can fail to be a symmetric-kernel law: \cite{BDF} contains no occurrence of ``negative
association'', ``negative dependence'' or ``repulsi'', and its only remark on the point is the
warning that the negative-dependence references ``mostly deal with symmetric
kernels\ldots whereas all of our examples involve nonsymmetric kernels''. This is a statement
about what those sources \emph{say}, not about what is decidable: for the two-block-factor
process displayed in~(b) the three-point NA inequalities can be tested directly from the atoms
recorded there, and we do not present that test as open. Our search could not
reach MathSciNet or Google Scholar (Section~\ref{sec:open}), so no priority is claimed here.
What is offered is the explicit witness with its exact verification, the impossibility at
$n\le2$, and the codimension count of Section~\ref{sec:generic}.

Finally, the direction of Theorem~\ref{thm:A} used here is the \emph{elementary} one
(symmetric $\Rightarrow p=q\Rightarrow\sigma^{2}=4uvw$). The hard converse -- when principal
minors determine a matrix up to diagonal similarity and transposition -- is the subject of
\cite{Loewy,Stevens,Mantelos,Arvind}, and is not needed.

\section{What is not settled}\label{sec:open}

\begin{itemize}
\item \textbf{Theorem~\ref{thm:generic} has no Hermitian analogue.} The law
$(\tfrac12,\tfrac12,\tfrac12;\tfrac1{100},\tfrac1{100},\tfrac1{100};0)$ has atoms
$\tfrac{11}{100},\tfrac{13}{100}(\times3),\tfrac{13}{100}(\times3),\tfrac{11}{100}$, is NA
with slack $1/100$, and \emph{is} the law of the Hermitian kernel
$\begin{psmallmatrix}1/2&1/10&-i/10\\1/10&1/2&1/10\\ i/10&1/10&1/2\end{psmallmatrix}$; all
inequalities being strict, an open $7$-dimensional set of NA laws is Hermitian-representable.
The Hermitian exclusion in Theorem~\ref{thm:main} is a \emph{pointwise} statement about $W$,
not a genericity statement.
\item \textbf{$\DPP(W)$ is not strongly Rayleigh}, so the analogous existence question for
strongly Rayleigh measures is \emph{not} answered here. With
$f(x)=\sum_{A}\mathbb P(X=A)\prod_{i\in A}x_{i}$ and
$D_{ij}f=\partial_{i}f\,\partial_{j}f-f\,\partial_{i}\partial_{j}f$, Br\"and\'en's criterion
for multiaffine polynomials \cite{Branden} requires $D_{ij}f\ge0$ on $\bbR^{3}$; at
$x=(-4,-4,-2)$ and $(i,j)=(1,2)$ one computes
$\partial_{1}f=\partial_{2}f=2735/10^{4}$, $f=-7675/10^{4}$,
$\partial_{1}\partial_{2}f=-987/10^{4}$, so
$D_{12}f=(2735^{2}-7675\cdot987)/10^{8}=-19/20000<0$. $\DPP(W)$ \emph{is} negatively
associated, but any
restatement of this result must say ``negatively associated'', never ``strongly Rayleigh''.
\item \textbf{The continuous-space version of the question is untouched.} What is settled here
is the finite-ground-set question of \cite{Poinas}; ``Continuous setting'' is a separate
subsubsection of the same discussion section.
\item \textbf{Prior-art coverage.} MathSciNet and Google Scholar were not reachable to us, so
a scoop in an unpublished thesis, in lecture notes, or in a work indexed only by MR cannot be
excluded.
\end{itemize}

\section{Reproduction and verification}\label{sec:verify}

Every object above is printed in full, in exact rationals, so the entire verification can be
redone with a pen: the eight atoms of $W$ by \eqref{eq:mobius}, the three slack classes, the
six-term determinant, Lemma~\ref{lem:pq}, and Theorem~\ref{thm:main}(iii) in four
multiplications and one subtraction. Remark~\ref{rem:signs} replaces even
Theorem~\ref{thm:A} by an eight-case enumeration.

The accompanying program \texttt{verify.py} (Python~3.9+, standard library only, exact
\texttt{Fraction} arithmetic, no float decision anywhere) reads the matrices \emph{as printed
above} and re-derives the numbers in the results of this paper, apart from the two groups
listed at the end of this section: the minors and atoms of $W$ by two
independent routes, negative association brute-forced \emph{from} \eqref{eq:na} over all
ordered pairs of disjoint nonempty sets and all pairs of proper nonempty up-closed families
with up-closure \emph{tested} rather than assumed, the invariants of \eqref{eq:inv} with
$u$ and $\sigma$ cross-checked against the principal minors, the eight-sign-pattern exclusion
of Remark~\ref{rem:signs}, the identity \eqref{eq:herm} on $1728$ Hermitian candidates with
the forced moduli, the padding and the $n\le2$ sweep of
Theorem~\ref{thm:D}, and the strong-Rayleigh point above. Its
up-set enumerator is pinned against the Dedekind numbers $2,3,6,20,168,7581$
(OEIS A000372) and its test counts against the prediction $30,348,4560,140058$ for
$n=3,4,5,6$. The recorded run reports
\texttt{VERDICT: ALL 71 CHECKS PASS} and its transcript, together with the scope statement the
program prints for itself, accompanies this paper.

\medskip
\noindent\textbf{Reading the transcript against this paper.} Its sections 1 and 2 verify
Theorem~\ref{thm:main}(i) and (ii); its section 3 verifies Lemma~\ref{lem:pq},
Theorem~\ref{thm:main}(iii) and Remark~\ref{rem:signs}; its section 8, headed ``Theorem D'',
verifies the padding and
$n\le2$ parts of Theorem~\ref{thm:D}; its section 10 verifies the strong-Rayleigh point of
Section~\ref{sec:open}. Its remaining headings name objects not discussed above and on which
nothing above depends: the controls of its section 4, the closed-form slack formula it calls
``Theorem B'', a second $n=3$ witness $W'$, the $\sigma$-interval it calls ``Theorem C'', the
$n=4$ witness, the box certificate of its section 9, and the conditional-association checks.
Those labels are the program's own and do not correspond to statements of this paper.

\medskip
\noindent\textbf{What the program does \emph{not} cover.} Two groups of numbers appearing
above are outside its $71$ checks, and are stated here so that a reader does not look for
them in the transcript. They are elementary and were checked by hand only.
\begin{itemize}
\item The arithmetic of Section~\ref{sec:attribution}: the triple invariants of Poinas'
coupling ($u=w=b^{2}$, $v=K_{11}(K_{11}-p_{1})$, gap $(b^{2}p_{1})^{2}$), and the
Borodin--Diaconis--Fulman numbers ($\det K=0$ against a symmetric competitor's $-d^{3}$, the
external checksum $\mathbb P(0,0,0)=1-3d+d^{2}$, and $u=w=1/12$, $v=0$ for the descent
process). The program's ``source's own kernel'' control is the uniform-on-even-subsets
kernel, a \emph{different} object from the two discussed there.
\item The Hermitian counterexample in the first item of Section~\ref{sec:open} (atoms
$11/100$ and $13/100$, NA slack $1/100$). The program verifies the Hermitian
\emph{inequality} \eqref{eq:herm} on $1728$ candidates, which is what
Theorem~\ref{thm:main}(iii) needs; it does not evaluate that particular kernel.
\end{itemize}
\noindent
Neither group is load-bearing for any theorem above: the first is attribution, the second is
a disclaimer that only weakens Theorem~\ref{thm:generic}.

\begin{thebibliography}{9}

\bibitem{Arvind}
V.~Arvind, A.~Chatterjee, P.~Mukhopadhyay, and others,
\emph{Characterizing and testing principal minor equivalence of matrices},
\href{https://arxiv.org/abs/2410.01961}{arXiv:2410.01961};
\href{https://doi.org/10.1145/3717823.3718146}{doi:10.1145/3717823.3718146}.

\bibitem{BBL}
J.~Borcea, P.~Br\"and\'en, and T.~M. Liggett,
\emph{Negative dependence and the geometry of polynomials},
J. Amer. Math. Soc. \textbf{22} (2009), no.~2, 521--567.
\href{https://doi.org/10.1090/S0894-0347-08-00618-8}{doi:10.1090/S0894-0347-08-00618-8}.

\bibitem{BDF}
A.~Borodin, P.~Diaconis, and J.~Fulman,
\emph{On adding a list of numbers (and other one-dependent determinantal processes)},
Bull. Amer. Math. Soc. \textbf{47} (2010), no.~4, 639--670.
\href{https://doi.org/10.1090/S0273-0979-2010-01306-9}{doi:10.1090/S0273-0979-2010-01306-9}.

\bibitem{Branden}
P.~Br\"and\'en,
\emph{Polynomials with the half-plane property and matroid theory},
Adv. Math. \textbf{216} (2007), no.~1, 302--320.
\href{https://doi.org/10.1016/j.aim.2007.05.011}{doi:10.1016/j.aim.2007.05.011}.

\bibitem{JDP}
K.~Joag-Dev and F.~Proschan,
\emph{Negative association of random variables with applications},
Ann. Statist. \textbf{11} (1983), no.~1, 286--295.
\href{https://doi.org/10.1214/aos/1176346079}{doi:10.1214/aos/1176346079}.

\bibitem{Loewy}
R.~Loewy,
\emph{Principal minors and diagonal similarity of matrices},
Linear Algebra Appl. \textbf{78} (1986), 23--64.
\href{https://doi.org/10.1016/0024-3795(86)90015-7}{doi:10.1016/0024-3795(86)90015-7}.

\bibitem{Lyons}
R.~Lyons,
\emph{Determinantal probability measures},
Publ. Math. Inst. Hautes \'Etudes Sci. \textbf{98} (2003), 167--212.
\href{https://doi.org/10.1007/s10240-003-0018-x}{doi:10.1007/s10240-003-0018-x}.

\bibitem{Mantelos}
H.~Sapranidis Mantelos,
\emph{Determinantally equivalent nonzero functions},
\href{https://arxiv.org/abs/2302.02471}{arXiv:2302.02471};
\href{https://doi.org/10.1016/j.disc.2026.115021}{doi:10.1016/j.disc.2026.115021}.

\bibitem{Poinas}
A.~Poinas,
\emph{On determinantal point processes with nonsymmetric kernels},
Electron. J. Probab. \textbf{31} (2026), no.~94.
\href{https://doi.org/10.1214/26-ejp1550}{doi:10.1214/26-ejp1550};
\href{https://arxiv.org/abs/2406.03360}{arXiv:2406.03360}.
The quotation in Section~\ref{sec:question} is from the \LaTeX{} e-print of \textbf{v4}, the
post-acceptance version of record.

\bibitem{Stevens}
M.~Stevens,
\emph{Equivalent symmetric kernels of determinantal point processes},
Random Matrices Theory Appl. \textbf{10} (2021), no.~3, 2150027.
\href{https://doi.org/10.1142/S2010326321500271}{doi:10.1142/S2010326321500271}.

\end{thebibliography}

\end{document}
